\documentclass[10pt,a4paper]{amsart}
\usepackage[margin=19mm]{geometry}
\usepackage{amsmath,amssymb,mathtools}
\usepackage[scr=rsfs,cal=euler]{mathalfa}
\usepackage{xfrac}
\usepackage[unicode,hidelinks,bookmarks=false]{hyperref}
\newcommand{\ie}{i.e.\ }
\newcommand{\cf}{cf.\ }
\newcommand{\resp}{respectively\ }
\newcommand{\Subsection}{\subsection}
\providecommand{\nobreakdash}{\nobreak\hbox{-}}
\setlength{\emergencystretch}{2em}
\multlinegap0pt
\usepackage{bm}
\usepackage{amssymb}
\usepackage{algorithm}\usepackage{algpseudocode}
\usepackage[shortlabels]{enumitem}
\newtheorem{definition}{Definition}
\newtheorem{lemma}{Lemma}
\newtheorem{theorem}{Theorem}
\title{Modeling AdaGrad, RMSProp, and Adam \\ with Integro-Differential Equations}
\author{Carlos Heredia}
\date{Source: arXiv:2411.09734v3 (CC BY 4.0)}
\begin{document}
\maketitle

\noindent\emph{Source and licence.} Modeling AdaGrad, RMSProp, and Adam \\ with Integro-Differential Equations. Version arXiv:2411.09734v3 (CC BY 4.0), distributed by arXiv under the Creative Commons Attribution 4.0 International licence, http://creativecommons.org/licenses/by/4.0/, as recorded in the arXiv licence metadata for this exact version and shown on the abstract page https://arxiv.org/abs/2411.09734v3. Full licence notice: \texttt{licenses/arxiv-2411.09734-CC-BY-4.0.txt}.\par\smallskip

\noindent\textit{Editorial scope.} Counted unit: the article's theorem thm:extended\_memory\_adagrad\_rmsprop (source0.tex lines 1569--1657). The article's complete original proof of it is delivered with it (source0.tex lines 1659--2030, one \texttt{proof} environment of 372 lines). No mathematics is written, repaired or re-derived: the statement and the proof are the article's own lines, in the article's own notation, and no retained line is substituted or re-flowed.  Retained uncounted as the source-local context the proof needs: lines 1078--1084; lines 1133--1141; lines 1214--1272.  Nothing was omitted from any retained span, so the package carries no omission record.  No retained line contains a citation, so the wrapper carries no bibliography and no bibliography entry.  No retained line contains a figure environment, an image-inclusion command or drawing code, so no image is delivered and the package declares no asset dependency.  Editorial formatting only, in addition: an AMS \texttt{amsart} preamble that restates the article's own declarations and macros used by the retained lines, namely the environments \texttt{definition}, \texttt{lemma}, \texttt{theorem} on separate counters, together with the standard packages those lines need.  The retained environments print in this document as Definition 1, Lemma 1, Theorem 1 in order of appearance, whereas the article prints the same items with the numbering of its own full document.  The complete native source of the article as distributed in the arXiv e-print for this exact version is bundled unchanged as \texttt{sources/168563/source0.tex}.
\begin{definition}\label{def:M}
Throughout, the notation \(M_\nu[g^2](t)\) is understood as
\[
    M_\nu[g^2](t) \ :=\ \int_0^t K_\nu(t-\tau)\,\|g(\tau)\|^2\,\mathrm{d}\tau.
\]
This convention avoids ambiguities that could arise from a component-wise interpretation.
\end{definition}

\begin{lemma}\label{lem:grad_bounds}
Let \( f : \mathbb{R}^n \to \mathbb{R} \) be \(m\)-strongly convex and \(L\)-smooth, with minimizer \(x^*\).  
Then, for all \( x \in \mathbb{R}^n \),
\[
    \frac{1}{2L} \, \|\nabla f(x)\|^2 
    \ \leq\ f(x) - f(x^*) 
    \ \leq\ \frac{1}{2m} \, \|\nabla f(x)\|^2.
\]
\end{lemma}

\begin{theorem}[Preliminary convergence bound for AdaGrad/RMSProp]
\label{thm:adagrad_rmsprop_memory}
Let \(f:\mathbb R^n\to\mathbb R\) be \(L\)-smooth and \(m\)-strongly convex, with unique minimizer \(\theta^*\). Let \(\nu\) be a finite positive Borel measure on \([0,\infty)\), and define \(M_\nu[g^2](t)\) as in Definition~\ref{def:M}. Assume that \(\theta(t)\) is a global classical solution of
\[
\dot\theta(t)
=
-
\frac{\nabla f(\theta(t))}
{\sqrt{M_\nu[g^2](t+\alpha)}+\epsilon},
\qquad t\ge 0,
\]
with \(\epsilon>0\) and \(\alpha>0\). Moreover, set
\[
\Phi_0:=f(\theta(0))-f(\theta^*),
\qquad
G_0:=L\sqrt{\frac{2\Phi_0}{m}},
\qquad
B:=G_0\sqrt{\|\nu\|}.
\]
Then, for every \(t\ge 0\),
\[
\|\theta(t)-\theta^*\|
\le
\|\theta(0)-\theta^*\|
\exp\!\left(-m\,I_\alpha(t)\right),
\]
where
\[
I_\alpha(t)
:=
\int_0^t
\frac{ds}{\epsilon+B\sqrt{s+\alpha}}.
\]
Equivalently, if \(B>0\),
\[
I_\alpha(t)
=
\frac{2}{B}
\left(\sqrt{t+\alpha}-\sqrt{\alpha}\right)
-
\frac{2\epsilon}{B^2}
\log\!\left(
\frac{\epsilon+B\sqrt{t+\alpha}}
{\epsilon+B\sqrt{\alpha}}
\right).
\]
In particular,
\[
I_\alpha(t)\sim \frac{2}{B}\sqrt t,
\qquad t\to\infty,
\]
and therefore
\begin{equation}\label{th1:bound}
\|\theta(t)-\theta^*\|
\le
C\exp(-c\sqrt t)
\end{equation}
for suitable constants \(C,c>0\).
\end{theorem}

\begin{theorem}[Extended memory functional for AdaGrad/RMSProp]
\label{thm:extended_memory_adagrad_rmsprop}
Let \(f:\mathbb R^n\to\mathbb R\) be \(L\)-smooth and \(m\)-strongly convex, with unique minimizer \(\theta^*\). Let \(g(t):=\nabla f(\theta(t))\), and consider the memory dynamics
\begin{equation}
\label{eq:shifted_memory_dynamics_thm2}
\dot\theta(t)
=
-
\frac{g(t)}
{\sqrt{M_\nu[g^2](t+\alpha)}+\epsilon},
\qquad t\ge 0,
\end{equation}
where \(\alpha>0\), \(\epsilon>0\), and \(M_\nu[g^2]\) is defined as in Definition~\ref{def:M}. Assume that the memory satisfies
\begin{equation}
\label{eq:memory_ode_thm2}
\dot M_\nu[g^2](t)
=
\lambda \|g(t)\|^2-\xi M_\nu[g^2](t),
\qquad 
\lambda>0,
\qquad 
\xi\in\{0,\lambda\}.
\end{equation}
Here \(\xi=0\) corresponds to AdaGrad and \(\xi=\lambda\) to RMSProp. Then there exists a constant \(\mu_0>0\) such that
\[
\frac{1}{\sqrt{M_\nu[g^2](t+\alpha)}+\epsilon}
\ge
\mu_0,
\qquad t\ge 0.
\]
Define
\[
\Phi(t):=f(\theta(t))-f(\theta^*),
\qquad
E_\rho(t):=\Phi(t)+\rho M_\nu[g^2](t),
\]
and choose
\[
\rho:=\frac{\mu_0}{2\lambda}.
\]
Then, for every \(t\ge 0\),
\begin{equation}
\label{eq:E_dissipation_thm2}
\dot E_\rho(t)
\le
-
m\mu_0\Phi(t)
-
\rho\xi M_\nu[g^2](t).
\end{equation}

Consequently:
\begin{enumerate}
\item If \(\xi=\lambda\), i.e. in the RMSProp case, then
\[
E_\rho(t)
\le
E_\rho(0)e^{-\kappa_0 t},
\qquad
\kappa_0:=\min\{m\mu_0,\xi\}>0.
\]
In particular,
\[
\Phi(t)\le E_\rho(0)e^{-\kappa_0 t},
\qquad
M_\nu[g^2](t)
\le
\frac{1}{\rho}E_\rho(0)e^{-\kappa_0 t},
\]
and
\[
\|\theta(t)-\theta^*\|
\le
\sqrt{\frac{2E_\rho(0)}{m}}\,e^{-\kappa_0 t/2}.
\]

\item If \(\xi=0\), i.e. in the AdaGrad case, then
\[
\dot E_\rho(t)\le -m\mu_0\Phi(t)\le 0.
\]
Hence \(E_\rho\) is nonincreasing. Moreover,
\[
\|\theta(t)-\theta^*\|
\le
\|\theta(0)-\theta^*\|e^{-m\mu_0 t},
\]
and \(M_\nu[g^2](t)\) converges to a finite limit as \(t\to\infty\). In general, this limit is not zero.
\end{enumerate}
\end{theorem}

\begin{proof}
We first treat the nontrivial case \(\theta(0)\neq \theta^*\). By Theorem~\ref{thm:adagrad_rmsprop_memory}, we know that
\[
\|\theta(t)-\theta^*\|
\le
R_0 e^{-mI_\alpha(t)},
\qquad
R_0:=\|\theta(0)-\theta^*\|.
\]
Since \(f\) is \(L\)-smooth and \(\nabla f(\theta^*)=0\), we have
\[
\|g(t)\|
=
\|\nabla f(\theta(t))-\nabla f(\theta^*)\|
\le
L\|\theta(t)-\theta^*\|.
\]
Therefore,
\begin{equation}
\label{eq:g_decay_from_thm1}
\|g(t)\|^2
\le
L^2R_0^2 e^{-2mI_\alpha(t)}.
\end{equation}

Moreover, by the explicit expression of \(I_\alpha(t)\), we have
\[
I_\alpha(t)\sim \frac{2}{B}\sqrt t,
\qquad t\to\infty.
\]
Hence
\[
e^{-2mI_\alpha(t)}
\sim
e^{-\frac{4m}{B}\sqrt t},
\qquad t\to\infty.
\]
Since \(\exp(-c\sqrt t)\) is integrable on \([0,\infty)\) for every \(c>0\), it follows that
\[
\mathcal J_\alpha
:=
\int_0^\infty e^{-2mI_\alpha(s)}\,\mathrm ds
<
\infty.
\]

We now prove that the memory is uniformly bounded.  Solving the linear equation
\begin{equation}\label{eq:M_nu_differential}
\dot M_\nu[g^2](t)
=
\lambda\|g(t)\|^2-\xi M_\nu[g^2](t)
\end{equation}
gives
\[
M_\nu[g^2](t)
=
e^{-\xi t}M_\nu[g^2](0)
+
\lambda
\int_0^t
e^{-\xi(t-s)}\|g(s)\|^2\,\mathrm ds.
\]
Since \(\xi\ge 0\), we have \(e^{-\xi(t-s)}\le 1\). Hence, using equation~(\ref{eq:g_decay_from_thm1}),
\[
M_\nu[g^2](t)
\le
M_\nu[g^2](0)
+
\lambda
\int_0^t
\|g(s)\|^2\,\mathrm ds
\le
M_\nu[g^2](0)
+
\lambda L^2R_0^2
\int_0^t
e^{-2mI_\alpha(s)}\,\mathrm ds.
\]
Therefore,
\[
M_\nu[g^2](t)
\le
M_\nu[g^2](0)
+
\lambda L^2R_0^2\mathcal J_\alpha
=
\overline M,
\qquad
t\ge 0.
\]
In particular,
\[
M_\nu[g^2](t+\alpha)\le \overline M,
\qquad
t\ge 0.
\]
Consequently,
\[
\sqrt{M_\nu[g^2](t+\alpha)}+\epsilon
\le
\sqrt{\overline M}+\epsilon,
\]
and so
\begin{equation}
\label{eq:mu_lower_thm2}
\mu(t)
:=
\frac{1}{\sqrt{M_\nu[g^2](t+\alpha)}+\epsilon}
\ge
\frac{1}{\epsilon+\sqrt{\overline M}}
=
\mu_0.
\end{equation}

We now compute the derivative of the objective gap $\Phi$. Since
\[
\dot\theta(t)
=
-
\mu(t)g(t),
\]
we obtain
\begin{equation}
\label{eq:Phi_dot_thm2}
\dot\Phi(t)
=
\langle g(t),\dot\theta(t)\rangle
=
-\mu(t)\|g(t)\|^2.
\end{equation}
For
\[
E_\rho(t)
=
\Phi(t)+\rho M_\nu[g^2](t),
\]
we have
\[
\dot E_\rho(t)
=
\dot\Phi(t)+\rho \dot M_\nu[g^2](t).
\]
Substituting equations (\ref{eq:M_nu_differential}) and (\ref{eq:Phi_dot_thm2}) gives
\begin{equation}
\label{eq:E_dot_raw_thm2}
\dot E_\rho(t)
=
-
\bigl(\mu(t)-\rho\lambda\bigr)\|g(t)\|^2
-
\rho\xi M_\nu[g^2](t).
\end{equation}
Choosing
\[
\rho
=
\frac{\mu_0}{2\lambda},
\]
and using \(\mu(t)\ge \mu_0\), we get
\[
\mu(t)-\rho\lambda
\ge
\mu_0-\frac{\mu_0}{2}
=
\frac{\mu_0}{2}.
\]
Therefore,
\begin{equation}
\label{eq:E_dot_first_bound_thm2}
\dot E_\rho(t)
\le
-
\frac{\mu_0}{2}\|g(t)\|^2
-
\rho\xi M_\nu[g^2](t).
\end{equation}

By Lemma~\ref{lem:grad_bounds}, since \(f\) is \(m\)-strongly convex and \(L\)-smooth,
\[
\Phi(t)
=
f(\theta(t))-f(\theta^*)
\le
\frac{1}{2m}\|g(t)\|^2.
\]
Hence
\[
\|g(t)\|^2
\ge
2m\Phi(t).
\]
Substituting this into equation~(\ref{eq:E_dot_first_bound_thm2}), we obtain
\[
\dot E_\rho(t)
\le
-
m\mu_0\Phi(t)
-
\rho\xi M_\nu[g^2](t),
\]
which proves equation~(\ref{eq:E_dissipation_thm2}). We now distinguish the two cases:

\medskip

\noindent\textbf{RMSProp case: \(\xi=\lambda\).}
If \(\xi=\lambda>0\), then
\[
\dot E_\rho(t)
\le
-
m\mu_0\Phi(t)
-
\rho\xi M_\nu[g^2](t).
\]
Let
\[
\kappa_0
:=
\min\{m\mu_0,\xi\}.
\]
Then
\[
m\mu_0\Phi(t)
+
\rho\xi M_\nu[g^2](t)
\ge
\kappa_0
\bigl(
\Phi(t)+\rho M_\nu[g^2](t)
\bigr)
=
\kappa_0 E_\rho(t).
\]
Therefore,
\[
\dot E_\rho(t)
\le
-\kappa_0 E_\rho(t).
\]
By Gronwall's inequality,
\[
E_\rho(t)
\le
E_\rho(0)e^{-\kappa_0 t}.
\]
Since both terms in \(E_\rho(t)=\Phi(t)+\rho M_\nu[g^2](t)\) are nonnegative, we get
\[
\Phi(t)
\le
E_\rho(0)e^{-\kappa_0 t},
\]
and
\[
M_\nu[g^2](t)
\le
\frac{1}{\rho}E_\rho(0)e^{-\kappa_0 t}.
\]
Finally, strong convexity gives
\[
\Phi(t)
\ge
\frac m2\|\theta(t)-\theta^*\|^2.
\]
Thus
\[
\|\theta(t)-\theta^*\|
\le
\sqrt{\frac{2}{m}\Phi(t)}
\le
\sqrt{\frac{2E_\rho(0)}{m}}e^{-\kappa_0 t/2}.
\]
This proves the RMSProp claims.

\medskip

\noindent\textbf{AdaGrad case: \(\xi=0\).}
If \(\xi=0\), then equation~(\ref{eq:E_dissipation_thm2}) becomes
\[
\dot E_\rho(t)
\le
-
m\mu_0\Phi(t)
\le 0.
\]
Hence \(E_\rho\) is nonincreasing. However, because there is no negative term proportional to \(M_\nu[g^2](t)\), this inequality cannot be closed in the form
\[
\dot E_\rho(t)\le -\kappa E_\rho(t)
\]
with some \(\kappa>0\). This reflects the purely accumulative character of AdaGrad. Nevertheless, the distance to the minimizer itself decays exponentially with the initial constant appearing in the statement. Indeed, define
\[
\mathcal L(t)
:=
\frac12\|\theta(t)-\theta^*\|^2.
\]
Since
\[
\dot\theta(t)=-\mu(t)g(t),
\qquad
\mu(t):=
\frac{1}{\sqrt{M_\nu[g^2](t+\alpha)}+\epsilon},
\]
we have
\[
\dot{\mathcal L}(t)
=
\langle \theta(t)-\theta^*,\dot\theta(t)\rangle
=
-\mu(t)\langle \theta(t)-\theta^*,g(t)\rangle.
\]
By \(m\)-strong convexity,
\[
\langle \theta(t)-\theta^*,g(t)\rangle
=
\langle \theta(t)-\theta^*,
\nabla f(\theta(t))-\nabla f(\theta^*)\rangle
\ge
m\|\theta(t)-\theta^*\|^2
=
2m\mathcal L(t).
\]
Using also \(\mu(t)\ge\mu_0\), we obtain
\[
\dot{\mathcal L}(t)
\le
-2m\mu_0\mathcal L(t).
\]
Therefore, by Gronwall's inequality,
\[
\mathcal L(t)
\le
\mathcal L(0)e^{-2m\mu_0t}.
\]
Taking square roots gives
\[
\|\theta(t)-\theta^*\|
\le
\|\theta(0)-\theta^*\|e^{-m\mu_0t}.
\]

Finally, in the AdaGrad case,
\[
\dot M_\nu[g^2](t)
=
\lambda\|g(t)\|^2
\ge 0,
\]
so \(M_\nu[g^2](t)\) is nondecreasing. Moreover, since \(f\) is \(L\)-smooth and \(\nabla f(\theta^*)=0\), the previous estimate gives
\[
\|g(t)\|
=
\|\nabla f(\theta(t))-\nabla f(\theta^*)\|
\le
L\|\theta(t)-\theta^*\|
\le
L\|\theta(0)-\theta^*\|e^{-m\mu_0t}.
\]
Hence \(\|g(t)\|^2\in L^1(0,\infty)\). Therefore
\[
M_\nu[g^2](t)
=
M_\nu[g^2](0)
+
\lambda\int_0^t\|g(s)\|^2\,\mathrm ds
\]
converges to a finite limit as \(t\to\infty\). In general, this limit is positive, so AdaGrad retains a persistent accumulated memory.

The case \(\theta(0)=\theta^*\) is immediate: then \(g(0)=0\), and the constant trajectory
\[
\theta(t)\equiv \theta^*
\]
solves the dynamics, with zero memory evolution if the memory is initialized at zero. This completes the proof.
\end{proof}

\end{document}
